\documentclass[11pt]{article}
\usepackage[margin=1in]{geometry}
\usepackage{amsmath,amssymb,amsthm}
\usepackage{hyperref}
\usepackage[numbers,sort&compress]{natbib}

\newtheorem{theorem}{Theorem}
\newtheorem{lemma}{Lemma}
\newtheorem{corollary}{Corollary}
\newtheorem{definition}{Definition}
\newtheorem{remark}{Remark}

\newcommand{\ste}[1]{\texttt{#1}}

\title{Adaptive Consistency: Enforcement, Completeness, and the Width
Formalisms\\
\large A literature review and annotated bibliography, oriented toward
the Lean mechanization in \ste{Ste.BucketConsistency}}
\author{Project \texttt{thejeb}}
\date{July 24, 2026}

\begin{document}
\maketitle

\begin{abstract}
The repository's \ste{Ste.AdaptiveConsistency} module mechanized the
\emph{sufficiency} half of the adaptive-consistency theorem: a bucket
network that is directionally consistent along an elimination order
admits a backtrack-free solution. This review surveys the literature
behind the \emph{other} half --- that adaptive consistency (bucket
elimination with existential projection) \emph{establishes} directional
consistency while \emph{preserving} the solution set, at cost
exponential only in the induced width --- and behind the identification
of induced width with graph treewidth. It then maps the surveyed
results onto the repository's existing Lean modules
(\ste{Ste.Elimination}, \ste{Ste.Projection}, \ste{Ste.Treedecomp},
\ste{Ste.GraphTreewidth}, \ste{Ste.AdaptiveConsistency}) and onto the
new module \ste{Ste.BucketConsistency}, stating precisely which
fragments are machine-checked and which remain open. Following the
project rule --- distrust results that are not machine proven --- all
claims about the mechanization are limited to statements that compile
in CI; everything else is attributed to the cited literature.
\end{abstract}

\section{The problem}

A finite constraint network consists of variables $x_0,\dots,x_n$ with
domains $D_i$ and a set of relations (constraints), each over a subset
of the variables (its \emph{scope}). Three questions dominate the
structural theory: (i)~when does \emph{local} consistency imply
\emph{global} solvability; (ii)~can local consistency of the required
strength be \emph{enforced} by an algorithm that does not change the
solution set; and (iii)~what \emph{graph parameter} of the constraint
network governs the cost of that enforcement. The classical answers ---
directional $(w{+}1)$-consistency suffices along an order of induced
width $w$ \citep{freuder1982backtrack, dechterpearl1987network};
adaptive consistency / bucket elimination enforces it in
$O(n\,k^{w+1})$ time and space while preserving the solution set
\citep{dechterpearl1987network, dechter1999bucket,
dechter2003constraint}; and induced width is graph treewidth
\citep{arnborg1987complexity, bodlaender1998arboretum} --- together
constitute the tractability theory of bounded-width constraint solving,
culminating in Grohe's converse \citep{grohe2007complexity}.

The repository proved (i) in Lean; this review is the survey component
of the attack on (ii) and (iii). Section~\ref{sec:classical} reviews
the classical results with an eye to what a mechanization must supply;
Section~\ref{sec:width} reviews the width formalisms;
Section~\ref{sec:mech} audits the mechanization, including a gap the
enforcement analysis surfaced in the repository's own statement of
directional consistency; Section~\ref{sec:bib} is the annotated
bibliography.

\section{Local consistency and backtrack-free search}
\label{sec:classical}

\subsection{From arc consistency to width}

Mackworth systematized the local-consistency propagation algorithms
(node, arc, and path consistency) as fixed-point computations on a
network of relations \citep{mackworth1977consistency}. Arc consistency
alone does not imply solvability; the decisive insight is Freuder's:
what matters is the \emph{relationship} between the level of local
consistency and a structural parameter of the constraint graph.
Freuder defined the \emph{width} of a vertex ordering as the maximal
number of earlier neighbours of any vertex, and the width of a graph
as the minimum over orderings, and proved: if a network is strongly
$(w{+}1)$-consistent and its constraint graph has width $w$, then a
solution can be assembled greedily along the ordering with no
backtracking \citep{freuder1982backtrack}. For trees ($w=1$) arc
consistency suffices --- the case mechanized in
\ste{Ste.ConsistencyTree}. Freuder later generalized the statement to
\emph{backtrack-bounded} search: $b$-level consistency relative to
orderings with bounded ``width at level $b$'' bounds the depth of
backtracking rather than eliminating it
\citep{freuder1985backtrackBounded}; the $k$-tree connection implicit
there --- graphs of width $w$ in Freuder's sense are exactly partial
$w$-trees --- is the bridge to the treewidth literature
(Section~\ref{sec:width}).

The subtlety that drives everything afterwards: enforcing
$(w{+}1)$-consistency by the naive closure algorithms \emph{adds
constraints}, which adds edges to the constraint graph, which can
\emph{increase its width}. Freuder's theorem therefore does not by
itself yield an algorithm: the hypothesis (consistency level) and the
parameter (width) chase each other. Two escapes were found. Seidel's
``invasion'' procedure processes the variables in one direction and
records, for each frontier, the set of partial solutions --- solving
the network in time exponential in the maximal frontier size
\citep{seidel1981invasion}. Dechter and Pearl's \emph{adaptive
consistency} is the refined form: rather than enforcing a uniform
consistency level everywhere, adapt the level per variable, exactly
matching each variable's set of earlier neighbours \emph{in the graph
as modified by the algorithm's own recorded constraints}
\citep{dechterpearl1987network}.

\subsection{Adaptive consistency and bucket elimination}

Fix an ordering $x_0 < x_1 < \dots < x_n$ and process variables from
last to first. At $x_i$, collect the \emph{bucket}: all constraints
(original or previously recorded) whose scope contains $x_i$ and no
later variable. Join them, and \emph{project out} $x_i$
existentially: the recorded constraint holds of an assignment to the
remaining scope variables iff \emph{some} value of $x_i$ extends it to
satisfy the whole bucket. Place the recorded constraint in the bucket
of its highest remaining variable and continue. This is adaptive
consistency \citep{dechterpearl1987network}, later renamed and
generalized as the constraint instance of \emph{bucket elimination},
Dechter's unifying dynamic-programming scheme that also captures
directional resolution, Fourier elimination, and the elimination
algorithms of probabilistic inference \citep{dechter1999bucket}. The
textbook treatment is \citet[ch.~4]{dechter2003constraint}.

Three claims make adaptive consistency a \emph{complete solver}, and a
mechanization must separate them cleanly:

\begin{enumerate}
\item \textbf{Solution-set preservation.} Each recorded constraint is
  a \emph{projection} of a join of existing constraints, hence implied
  by them: adding it changes nothing for full assignments, and
  restricting to the remaining variables loses exactly the eliminated
  coordinate. Formally, if $S$ is the solution set of the current
  network over variables $V$ and the step eliminates $v$, the new
  network's solution set over $V$ is $\exists_v S$ --- the projection.
  In particular feasibility is preserved \emph{exactly}, with no
  substituted values; this is where projection differs essentially
  from the conditioning (value-substitution) step, whose verdict is
  conditional on the chosen values.
\item \textbf{Enforcement of directional consistency.} After the run,
  every assignment to $x_0,\dots,x_{i-1}$ that satisfies all
  constraints among them --- \emph{including the recorded ones} ---
  extends to $x_i$: the recorded projection of bucket $i$ is exactly
  the condition for a witness to exist, and it was placed among the
  earlier constraints. Note the logical form: the extension property
  is \emph{relative} to the earlier constraints, not absolute; this
  distinction matters for the mechanization
  (Section~\ref{sec:gap}).
\item \textbf{Cost.} Each bucket's join and projection touch only the
  variables in that bucket's scopes; if the ordering has induced width
  $w$ (Section~\ref{sec:width}), every recorded scope has at most $w$
  variables, so each table has at most $k^{w+1}$ rows over domains of
  size $k$, and the total work is $O(n\,k^{w+1})$
  \citep{dechterpearl1987network, dechter2003constraint}.
\end{enumerate}

Given (1)--(3), the backtrack-free construction of
\citet{freuder1982backtrack} (in its directional form) extracts a
solution with no search, and the decision ``is the network
solvable?''\ is read off the final state. Dechter and Pearl also
proved the ordering-graph identification that makes (3) sharp, and
tree clustering --- the join-tree compilation whose cluster sizes are
governed by the same parameter --- as the dual, query-oriented
formulation \citep{dechterpearl1989tree}.

\section{The width formalisms}
\label{sec:width}

Three formalisms measure the same quantity.

\emph{Induced width (elimination width).} For an ordering $\sigma$ of
the vertices of a graph $G$, process vertices from last to first,
connecting the earlier neighbours of each processed vertex (the
\emph{fill-in}); the induced width $w^*(\sigma)$ is the maximal number
of earlier neighbours in the filled graph, and the induced width of
$G$ is the minimum over orderings. This is precisely the maximal
recorded-scope size of adaptive consistency along $\sigma$
\citep{dechterpearl1987network, dechter2003constraint}, and it is the
``dimension'' of Bertel\`e and Brioschi's nonserial dynamic
programming, where the notion (under the name of that book) predates
the constraint literature by fifteen years \citep{bertele1972nonserial}.

\emph{Partial $k$-trees.} A $k$-tree is built from a $(k{+}1)$-clique
by repeatedly attaching a new vertex adjacent to a $k$-clique; partial
$k$-trees are their subgraphs. Arnborg, Corneil and Proskurowski
proved that recognizing partial $k$-trees --- equivalently, deciding
treewidth --- is NP-hard when $k$ is part of the input, and gave an
$O(n^{k+2})$ algorithm for fixed $k$ \citep{arnborg1987complexity}.
This is why every completeness theorem quantifies over a \emph{given}
ordering: finding the optimal one is itself intractable.

\emph{Tree decompositions.} Robertson and Seymour introduced
tree decompositions and treewidth in the Graph Minors series and
showed the parameter's algorithmic significance
\citep{robertson1986graphminorsII}. Bodlaender's survey
\citep{bodlaender1998arboretum} assembles the equivalences: a graph
has treewidth $\le w$ iff it is a partial $w$-tree iff it has an
elimination ordering of induced width $\le w$ --- the elimination
ordering being read off a tree decomposition by eliminating leaves of
bags, and conversely the fill-in cliques of an ordering forming the
bags of a decomposition. It is exactly this last construction that
\ste{Ste.GraphTreewidth} mechanizes from the bucket-elimination side.

The circle closes with the converse direction: Grohe proved (for
bounded arity, under a standard parameterized-complexity assumption)
that bounded treewidth of the core is the \emph{only} structural reason
constraint classes are tractable \citep{grohe2007complexity} ---
adaptive consistency at bounded width is not just \emph{a} tractable
algorithm but essentially \emph{the} tractable frontier.

\section{The mechanization: what is proved, and where}
\label{sec:mech}

\subsection{Prior state of the repository}

\begin{itemize}
\item \ste{Ste.Elimination}: bucket elimination with a
  \emph{substitutive} step (\ste{bucketHead} joins the bucket of $v$
  and conditions on a chosen value $v \mapsto a$).
  Machine-proved: the fold faithfully computes the elimination of the
  joint constraint (\ste{joinConstraint\_bucketEliminate}), a complete
  order decides the \emph{substituted} instance
  (\ste{bucketEliminate\_decides}), and the total table space is
  $n\,k^{w+1}$ given bags of size $\le w{+}1$
  (\ste{bucketEliminate\_total\_space}). The verdict is conditional on
  the substituted values --- claim (1) of
  Section~\ref{sec:classical} fails for this step.
\item \ste{Ste.Projection}: the existential step \ste{project} on
  monolithic constraints. Machine-proved: projection preserves
  nonemptiness exactly (\ste{project\_nonempty\_iff}) and a complete
  projective run decides feasibility unconditionally
  (\ste{projectEliminate\_decides\_feasibility}) --- but with no bucket
  data structure, hence no width accounting.
\item \ste{Ste.Treedecomp}: \ste{AchievesWidth} and the attained
  minimum \ste{inducedTreewidth}, discharging the free width
  hypothesis of the space bound.
\item \ste{Ste.GraphTreewidth}: tree decompositions from scratch, the
  elimination decomposition, and the headline
  \ste{treewidth\_primalGraph\_le}:
  $\mathrm{tw}(\text{primal}) \le \ste{inducedTreewidth} + 1$.
\item \ste{Ste.AdaptiveConsistency}: the bucket-network formalism
  ($\ste{sep}\ i$, $\ste{R}\ i$, \ste{SepSupported},
  \ste{SepWidthLE}) and the sufficiency theorem
  \ste{directionalConsistent\_solvable}: nonempty domains +
  directional consistency $\Rightarrow$ a global solution, greedily.
\end{itemize}

\subsection{A gap surfaced by the enforcement analysis}
\label{sec:gap}

\ste{Ste.AdaptiveConsistency}'s \ste{DirectionalConsistent} demands,
for \emph{every} domain-respecting assignment $g$, a witness $c \in
D_{i+1}$ with $R_i(g[x_{i+1} \mapsto c])$. But what bucket elimination
establishes at bucket $i$ is weaker: a witness exists for assignments
that satisfy the \emph{earlier buckets} --- because the recorded
projection of bucket $i$ is itself one of those earlier constraints.
Dechter's directional consistency is the relative notion
\citep[ch.~4]{dechter2003constraint}. Thus the repository's
sufficiency theorem, as previously stated, consumed a hypothesis
strictly stronger than what the enforcement algorithm can supply: the
two halves did not compose. This observation --- easy to miss on paper,
where ``consistent partial assignment'' silently carries the recorded
constraints --- is, we believe, the kind of precision dividend one
mechanizes for.

\subsection{The new module \ste{Ste.BucketConsistency}}

All statements below compile in CI with no \ste{sorry}, \ste{admit},
or added axioms; the run number is recorded in the accompanying log
entry.

\textbf{Part 1 (enforcement at the joint-constraint level).} The
projective bucket step \ste{projectBucketStep} replaces the bucket of
$v$ by its joined \emph{projection} (the recorded message is
definitionally \ste{project} of the joined bucket,
\ste{projectBucketHead\_snd}). Machine-proved:
\begin{itemize}
\item \ste{joinConstraint\_projectBucketStep}: the residual state's
  joint constraint \emph{equals} $\exists_v(\text{previous joint
  constraint})$ --- claim (1), solution-set preservation, exactly;
\item \ste{exists\_update\_mem\_of\_mem\_projectBucketStep}: every
  joint solution of the residual state extends at $v$ alone to a joint
  solution of the previous state --- claim (2) in instance form, the
  extension property the recorded projection buys;
\item \ste{joinConstraint\_projectBucketEliminate},
  \ste{projectBucketEliminate\_nonempty\_iff},
  \ste{projectBucketEliminate\_decides}: the fold computes
  \ste{projectEliminate}, preserves feasibility exactly, and a
  complete order yields the \emph{unconditional} feasibility verdict
  --- projective bucket elimination is a complete solver;
\item \ste{exists\_extension\_of\_mem\_projectBucketEliminate}:
  backtrack-free extraction --- any survivor of the full run extends
  to a solution of the original instance changing only eliminated
  variables.
\end{itemize}

\textbf{Part 2 (the relative sufficiency theorem).} The gap of
Section~\ref{sec:gap} is closed from the sufficiency side:
\ste{BucketRelativeConsistent} formalizes Dechter's relative
directional consistency, and
\ste{bucketRelativeConsistent\_solvable(\_from)} proves it implies
global solvability. The proof is a \emph{reduction}, not a new greedy
induction: gate each bucket by its predecessors (\ste{guardedR}),
observe that the gated network has the same solution set
(\ste{adaptiveSolution\_guardedR\_iff}, by strong induction on the
node index), is \ste{SepSupported} for the fat prefix separators, and
is \emph{unconditionally} directionally consistent (an update at
$x_{i+1}$ is invisible to earlier buckets, so an unsatisfied guard
stays unsatisfied), then invoke the existing engine. Solvability needs
no width bound, so the fat separators are harmless.

\textbf{Part 3 (the width link).} Machine-proved: the projective and
substitutive folds materialize identical scope traces
(\ste{projectBucketBags\_map\_fst} --- scopes evolve by scope-only
rules), so \ste{Ste.Treedecomp}'s width accounting transfers verbatim
(\ste{projectBucketBags\_card\_le}). Along the decreasing order on
$\mathrm{Fin}(n{+}1)$ --- the order implicit in
\ste{Ste.AdaptiveConsistency}'s numbering --- the recorded message
scope at the step eliminating $x_{i+1}$ defines \ste{runSep}, with:
\ste{runSep\_eq\_messageScope} (it is literally the
\ste{bucketHead}-style scope of the reached state),
\ste{runSep\_precedes} (topological order: the \ste{hsep} hypothesis
of the solvability theorems, now \emph{established} by the run), and
\ste{runSep\_widthLE} / \ste{runSep\_link}: if the decreasing order
witnesses $\ste{AchievesWidth}\ B\ w$ then
$\ste{SepWidthLE}\ (\ste{runSep}\ B)\ (w{+}1)$ and
$\ste{inducedTreewidth}\ B \le w$. Via
\ste{treewidth\_primalGraph\_le}, the separator bound of the
bucket-network formalism, the elimination width, and (up to the fixed
offset) primal-graph treewidth are controlled by the same run.

\subsection{Honest accounting of what remains open}
\label{sec:open}

\begin{enumerate}
\item \textbf{The per-node reassembly.} Part 1 preserves solutions and
  extracts them at the \emph{joint-constraint} level; Part 2 proves
  the \emph{per-node} relative condition sufficient. The missing link
  is the reassembly: packaging the final state of a run over
  $\mathrm{Fin}(n{+}1)$ as a literal bucket network $(D, \ste{runSep},
  R)$ with $\ste{adaptiveSet}\ D\ R = \ste{joinConstraint}\ B$ and
  \ste{BucketRelativeConsistent} $D$ $R$ --- requiring
  constraint-to-bucket provenance (each constraint joins the bucket of
  its maximal variable) and root-domain extraction (surviving
  scope-$\{x_0\}$ constraints become $D_0$). On paper this is
  bookkeeping; in Lean it is a substantial indexing argument, and it
  is \emph{not} mechanized. The completeness claim is therefore
  machine-proved in its joint form and paper-joined in its per-node
  form.
\item \textbf{General orders.} The width link is proved for the
  decreasing order; an arbitrary order is handled classically by
  renumbering the variables along it. The permutation transport is not
  mechanized.
\item \textbf{Time complexity.} Only table-\emph{space} bounds exist
  in the repository (\ste{Ste.Treewidth}, \ste{Ste.Elimination});
  running time is not modelled. The claim ``cost exponential only in
  $w$'' is machine-proved for space, cited for time
  \citep{dechterpearl1987network, dechter2003constraint}.
\end{enumerate}

\section{Annotated bibliography}
\label{sec:bib}

Each annotation is our own summary, written against the source's role
in this project.

\begin{description}

\item[\citet{freuder1982backtrack}] Defines the width of an ordered
constraint graph as the maximal number of earlier neighbours and
proves that strong $(w{+}1)$-consistency plus width $w$ yields
backtrack-free search; trees with arc consistency are the $w{=}1$
case. For this project it is the ancestral form of the sufficiency
half: \ste{Ste.ConsistencyTree} mechanizes its tree case and
\ste{Ste.AdaptiveConsistency} its order-directional generalization.
The paper's own limitation --- enforcing the required consistency can
raise the width --- is precisely what adaptive consistency later
repairs, and why enforcement is the harder half to state correctly.

\item[\citet{freuder1985backtrackBounded}] Extends the 1982 theorem
from backtrack-\emph{free} to backtrack-\emph{bounded}: $b$-level
relative consistency along orderings whose width is controlled at
level $b$ bounds the backtrack depth. The analysis runs through
$k$-trees, making explicit that Freuder's width-$w$ graphs are partial
$w$-trees --- the earliest bridge between the CSP width and the
graph-minors treewidth formalisms, and the reason
\ste{Ste.GraphTreewidth}'s bridge theorem has classical warrant.

\item[\citet{dechterpearl1987network}] Introduces adaptive consistency
--- per-variable directional consistency enforced by recording the
projected bucket constraint --- together with the induced-width
parameter $w^*$ and the $O(n\,k^{w^*+1})$ bound, and proves the
solution set is preserved because every recorded constraint is implied
by the ones it is derived from. This is the primary source for the
enforcement half attacked here: our
\ste{joinConstraint\_projectBucketStep} is its preservation lemma made
exact (equality with the projection, not merely implication), and our
\ste{BucketRelativeConsistent} is its output condition stated with the
relativization made explicit.

\item[\citet{dechterpearl1989tree}] Recasts the same machinery in
compiled form: build a join tree whose clusters are the maximal
cliques of a triangulation of the constraint graph, solve each
cluster, and propagate. Cluster size is induced width plus one, so
tree clustering and adaptive consistency have the same exponent; the
join-tree view is the one \ste{Ste.JunctionTree} and
\ste{Ste.Treedecomp} echo. Included because the equivalence of the
per-order (elimination) and per-decomposition (cluster) presentations
is exactly the width link mechanized in Part 3.

\item[\citet{dechter1999bucket}] Presents bucket elimination as a
single elimination scheme instantiated by a ``combine'' and a
``project'' operator, with adaptive consistency the
$\{\text{join},\exists\}$ instance and directional resolution,
Gaussian/Fourier elimination and probabilistic elimination as
siblings. Its clean separation of the data structure (buckets, one
pass) from the algebra (the two operators) is the design our
\ste{projectBucketHead}/\ste{projectBucketStep} follows: the module
swaps \ste{Ste.Elimination}'s substitutive operator for the
existential one and re-proves the fold correctness at that generality.

\item[\citet{dechter2003constraint}] The textbook consolidation
(ch.~4 for adaptive consistency and induced width). Valuable here for
its careful statement of \emph{directional} consistency as extension
of consistent partial assignments --- consistent \emph{with respect to
all constraints among the earlier variables, recorded ones included}
--- which is the reading our \ste{BucketRelativeConsistent} formalizes
and which the repository's earlier \ste{DirectionalConsistent}
unintentionally strengthened (Section~\ref{sec:gap}).

\item[\citet{mackworth1977consistency}] The systematization of node,
arc and path consistency and their propagation algorithms as
constraint-network fixed points. Included as the baseline the width
theory refines: uniform local consistency neither implies solvability
nor respects structure, and the entire adaptive line can be read as
replacing Mackworth's uniform closure with an order-adapted one.

\item[\citet{bertele1972nonserial}] The nonserial dynamic programming
monograph: eliminating variables from an objective decomposed over a
hypergraph, with cost governed by what is now called induced width
(their ``dimension''), including the NP-hardness flavour of choosing
the elimination order. Adaptive consistency is its Boolean
(feasibility) specialization; we cite it as the origin of the
elimination-width parameter that \ste{Ste.Treedecomp}'s
\ste{inducedTreewidth} mechanizes.

\item[\citet{seidel1981invasion}] The invasion procedure: sweep the
variables once, maintaining the set of partial solutions restricted to
the current frontier; time exponential in the maximal frontier.
Historically the first CSP decision procedure exponential only in a
separator-like parameter, and conceptually the ancestor of our
Part~1: the joint-constraint-level fold with exact feasibility
preservation is an invasion-style argument, with projections replacing
frontier tables.

\item[\citet{arnborg1987complexity}] Proves NP-completeness of
deciding treewidth (partial $k$-tree recognition) for variable $k$ and
gives the $O(n^{k+2})$ algorithm for fixed $k$. For this project it
justifies the shape of every mechanized statement: all width bounds
are stated for a \emph{given} order or decomposition, and
\ste{inducedTreewidth} is defined as an attained infimum rather than
computed.

\item[\citet{robertson1986graphminorsII}] Introduces tree
decompositions and treewidth and initiates their algorithmic use in
the Graph Minors project. The three axioms of
\ste{Ste.GraphTreewidth}'s \ste{TreeDecomposition} (vertex cover, edge
cover, running intersection in rooted form) are this paper's
definition, specialized to a list-of-bags presentation suitable for
Lean induction.

\item[\citet{bodlaender1998arboretum}] The survey establishing the
equivalence web around treewidth: partial $k$-trees, elimination
orderings, cops-and-robber games, brambles. The
elimination-order-to-decomposition direction surveyed there is the one
\ste{Ste.GraphTreewidth} mechanizes
(\ste{treewidth\_primalGraph\_le}); the converse direction
(decomposition to ordering, needed to replace ``the decreasing
order'' by ``some optimal order'' in Part 3 without renumbering) is
surveyed there too and remains unmechanized.

\item[\citet{grohe2007complexity}] The converse: for bounded-arity
classes, tractability of the homomorphism/CSP problem forces bounded
treewidth of the cores (assuming $\mathrm{FPT}\ne\mathrm{W[1]}$).
Cited to position the mechanized material: the enforcement theorems
formalize the \emph{algorithmic} side of a boundary that is, by this
result, information-theoretically tight.

\end{description}

\bibliographystyle{plainnat}
\bibliography{../../refs}

\end{document}
